% Part: proof-theory
% Chapter: cut-elimination
% Section: introduction

\documentclass[../../../include/open-logic-section]{subfiles}

\begin{document}

\olfileid{pt}{cut}{int}

\olsection{Pambuka}

Aturan \Cut{} yaiku
\[
\Axiom$\Gamma \fCenter \Delta, !A$
\Axiom$!A, \Pi \fCenter \Lambda$
\RightLabel{\Cut}
\BinaryInf$\Gamma, \Pi \fCenter \Delta, \Lambda$
\DisplayProof
\]
(!!{formula}~$!A$ diarani \emph{formula cut}.)

\Cut{} kalebu ing kalkulus sekèn Gentzen kang asli,~\Log{LK}.
\emph{Hauptsatz} (``teorema utama'') Gentzen mratélakaké yèn
saben sekèn kang !!{provable} ing~\Log{LK} uga
!!{provable} tanpa nganggo aturan \Cut{}; tegesé,
\Cut{} bisa ``diilangi'' saka bukti~\Log{LK}.
Sistem kita, \Log{G1c} lan \Log{G3c}, ora ngemot
\Cut, nanging pitakonan kang padha bisa diajokaké:
apa \Cut{} bisa diilangi saka !!{proof}s kang nganggo
\Cut{} bebarengan karo aturan liyané saka
\Log{G1c} (utawa saka \Log{G3c})?
Miturut istilah kang wis ditetepaké, pitakonan iki
padha karo: apa \Cut{} aturan admisibel ing
\Log{G1c} (utawa \Log{G3c})?

Teorema eliminasi cut mbokmenawa asil teori bukti kang
paling wigati lan misuwur. Sepisan, teorema iki nuduhaké
yèn \Log{G1c} lan \Log{G3c} ora mbutuhaké aturan kang
sakawit katon mesthi perlu. Nalika mbakuaké bukti nyata,
kita perlu cara kanggo migunakaké lema. Matématikawan
biasané mbuktèkaké asil, banjur migunakaké asil mau
kanggo mbuktèkaké asil liyané. Upamané, bukti asil $L$
(lema) dibakuaké minangka !!{proof} sekèn
$\Sequent L$. Banjur bukti asil kapindho~$R$ kang
nganggo~$L$ dibakuaké minangka !!{proof} sekèn
$L \Sequent R$. Kepriyé kita ngerti ana bukti kanggo~$R$,
yaiku !!a{proof} sekèn $\Sequent R$ waé?
Kita kudu bisa nggabungaké rong !!{proof}s iku,
lan aturan \Cut{} maringi cara kanggo nindakaké.

Kita wis nyebut yèn \Log{G1c} lan \Log{G3c} iku komplet,
mula bisa mbuktèkaké saben sekèn kang mesthiné bisa
dibuktèkaké, yaiku kabèh sekèn valid. Iki bisa
dibuktèkaké kanthi langsung. Nanging ing jaman Gentzen,
mung sistem derivasi aksiomatik kang dimangertèni komplet.
Kanggo nuduhaké yèn \Log{LK} komplet, Gentzen maringi
terjemahan !!{derivation}s saka sistem aksiomatik dadi
\Log{LK}-!!{proof}s. Terjemahan iki gumantung banget
marang \Cut{}. Eliminasi cut ora mung wigati kanggo
logika klasik: akèh logika liya nduwèni kalkulus sekèn.
Kanggo sawatara logika, bukti langsung kakompletan
kalkulus sekèné angel utawa ora mungkin, mula mung
terjemahan saka sistem liya (nganggo \Cut) kang bisa
nuduhaké kakompletan. Ing kahanan mangkono, bukti
eliminasi cut kanthi cara teori bukti dibutuhaké
kanggo nuduhaké yèn \Cut{} ora perlu lan sistem tanpa
\Cut{} iku komplet.

Eliminasi cut uga wigati amarga bisa digunakaké kanggo
ngolèhké asil liya kang menarik dhéwé. Rong tuladha kang
bakal kita rembug yaiku lema interpolasi Craig lan
teorema Herbrand.

Gagasan kanggo mbuktèkaké eliminasi cut lumrahé nuduhaké
carané !!a{proof} kang nganggo \Cut{} bisa diowahi dadi
bukti tanpa cut. Owah-owahan iki biasané ``lokal'':
kita njupuk pérangan saka !!a{proof} (lumrahé
sub-!!{proof} kang pungkasané nganggo inferènsi \Cut{})
lan nggantèni nganggo sub-!!{proof} liya kanggo sekèn
kang padha nanging inferènsi \Cut{}-é wis diilangi
utawa diganti inferènsi liya. Argumèn mangkono maringi
algoritma langkah demi langkah kanggo ngowahi !!{proof}s
nganggo \Cut{} dadi bukti tanpa cut. Nalika prosedur
eliminasi cut digunakaké, iki maringi informasi luwih akèh.
Tuladhané, bukti teori modèl kanggo lema interpolasi
bisa nuduhaké yèn ``yèn $!A \lif !B$'' valid, banjur
ana interpolan (kanggo saiki ora perlu dijlentrehaké
apa interpolan iku). Argumèn teori bukti kang nganggo
eliminasi cut maringi algoritma kang nampa !!a{proof}
kanggo $!A \lif !B$ lan ngasilaké interpolan.
Béda karo argumèn teori modèl, argumèn teori bukti
iki \emph{konstruktif}.

Ana rong cara lumrah kanggo mbuktèkaké teorema eliminasi
cut. Cara kapisan (wiwit Gentzen) nuduhaké yèn cut paling
dhuwur ing bukti bisa dibusak, banjur dibusak siji-siji
nganti entèk. Cara kapindho (wiwitané saka Sch\"utte lan Tait)
nggatekaké cut kang paling rumit ing !!a{proof}, banjur
mbusak cut maksimal siji-siji, yaiku ora ana cut liya
kang !!{formula} cuté luwih jero. Saben cara nduwèni
kaluwihan lan kakurangan. Ing sawatara sistem, cara siji
bisa lumaku nanging cara liyané angel utawa malah ora
mungkin. Mula kita njlentrehaké kaloro cara.
Kita bakal mbuktèkaké eliminasi cut kanggo~\Log{G3c}.
Satemah kang bakal kita tuduhaké bisa diilangi dudu
\Cut{} langsung, nanging \emph{cut kanthi kontèks padha}~\CutCS:
\[
\Axiom$\Gamma \fCenter \Delta, !A$
\Axiom$!A, \Gamma \fCenter \Delta$
\RightLabel{\CutCS}
\BinaryInf$\Gamma \fCenter \Delta$
\DisplayProof
\]
Yèn kalkulus sekèn ngidini pelemahan lan kontraksi
(minangka aturan nyata kaya ing \Log{G1c}, utawa
minangka aturan admisibel kaya ing \Log{G3c}),
\Cut{} bisa nyimulasi \CutCS, lan kosok balènané.
Kita milih \CutCS{} tinimbang \Cut amarga strategi
eliminasi cut kapisan luwih gampang dijlentrehaké kanggo
\CutCS, déné strategi kapindho mung lumaku kanggo~\CutCS.

\end{document}
